The \pvs library export required three separate developments:

Firstly, Sam Owre has extended \pvs with an XML export.  This is similar to the {\LaTeX}
extension in PVS, which is built on the Common Lisp Pretty Printing facility.  The XML
export was developed in parallel with a Relax NG specification for the PVS XML files.
Because PVS allows overloading of names, infers theory parameters, and automatically adds
conversions, the XML generation is driven from the internal type-checked abstract syntax,
rather than the parse tree.  Thus the generated XML contains the fully type-checked form
of a PVS specification with all overloading disambiguated.  Future work on this will
include the generation of XML forms for the proof trees.

\autoref{fig:pvsxml} exemplary shows the (slightly simplified) XML representation of the PVS theory presented in \autoref{fig:intro:prel:pvs:equivclos}.

\begin{figure}[ht]\centering
\begin{code}
 <theory place="6049 0 6075 22">
  <id>EquivalenceClosure</id>
  <const-decl place="6057 2 6058 75">
    <id>EquivClos</id>
    <arg-formals>
      <binding place="6057 12 6057 13">
        <id>R</id>
        <type-name>
          <id>PRED</id>
          <actuals>
            <tuple-type>
              <type-name><id>T</id></type-name>
              <type-name><id>T</id></type-name>
            </tuple-type>
          </actuals>
        </type-name>
      </binding>
    </arg-formals>
    <type-name place="6057 17 6057 31">
      <id>equivalence</id>
      <actuals>
        ...
\end{code}
\caption{A Simplified Part of the Function EquivalenceClosure/EquivClos in XML \protect\footnotemark}\label{fig:pvsxml}
\end{figure}
\footnotetext{\url{https://gl.mathhub.info/PVS/Prelude/blob/master/source/prelude/pvsxml/EquivalenceClosure.xml}}

Secondly, Florian Rabe documented the XML schema used by PVS as a set of case classes in Scala and wrote a generic XML parser in Scala that
generates a schema-specific parser from such a set of inductive types (see
\autoref{fig:pvsscala} for part of the specification).  That way any change to the
inductive types automatically changes the parser.  While seemingly a minor implementation
detail, this was critical for feasibility because the XML schema changed frequently along
the way.

\begin{figure}[ht]\centering
\begin{scalacode}
case class const_decl(
    named: ChainedDecl,
    arg_formals: List[bindings],
    tp: DeclaredType, 
    _def: Option[Expr]
) extends Decl
\end{scalacode}
\caption{The Scala-Specification of PVS Constant Declarations for XML Parsing}\label{fig:pvsscala}
\end{figure}

Thirdly, I wrote an \mmt plugin that parses the XML files generated by PVS and systematically translates their content into \ommt, using the symbols from my formalization of the PVS logic.  This includes creating various generic indexes that can be used later
for searching the content.

All processing steps preserve source references, i.e., URLs that point to a location (file
and line/column) in a source file (the quatruples of numbers at \lstinline|place=| in \autoref{fig:pvsxml} and
\lstinline|<link rel="...?sourceRef"| in \autoref{fig:pvsomdoc}).

\begin{figure}[ht]\centering
\begin{code}[basicstyle=\sf\scriptsize]
<omdoc>
 <theory name="EquivalenceClosure"
  base="http://pvs.csl.sri.com/prelude"
  meta="http://pvs.csl.sri.com/?PVS">
  <constant name="EquivClos">
   <type>
    <om:OMOBJ>
     <om:OMA>
      <om:OMS base="http://cds.omdoc.org/urtheories" module="LambdaPi" name="apply"/>
      <om:OMS base="http://pvs.csl.sri.com/" module="PVS" name="expr"/>
      <om:OMA>
       <om:OMS base="http://cds.omdoc.org/urtheories" module="LambdaPi" name="apply"/>
       <om:OMS base="http://pvs.csl.sri.com/" module="PVS" name="pvspi"/>
        ...
        <metadata>
         <link rel="http://cds.omdoc.org/mmt?metadata?sourceRef"
           resource="prelude/pvsxml/EquivalenceClosure.xml#-1.6057.17:-1.6057.31"/>
        </metadata>
      </om:OMA>
\end{code}
\caption{A Part of the Function EquivalenceClosure/EquivClos in OMDoc \protect\footnotemark}\label{fig:pvsomdoc}
\end{figure}
\footnotetext{\url{https://gl.mathhub.info/PVS/Prelude/blob/master/content/http..pvs.csl.sri.com/prelude/$Equivalence$Closure.omdoc}}

  The table in \autoref{fig:sizes} gives an overview of the sizes of the involved
  libraries and the run times\footnote{all numbers measured on standard laptops} of the
  conversion steps. We note that the XML encoding considerably increases the size of
  representations. This is due to two effects: the internal, disambiguated form contains
  significantly more information than the user syntax (e.g. theory parameter instances and
  reconstructed types), and XML as a machine-oriented format is naturally more
  verbose. Furthermore, \omdoc uses OpenMath for term structures, which again increases
  file size. While in practice the file sizes are no problem for the \mmt tools presented here,
  the analogous import for the Isabelle libraries due to Makarius Wenzel\footnote{As of yet unpublished} has
  prompted us to compress the \omdoc files by default, significantly reducing their sizes. The \mmt system can read
  the compressed files directly without need for a user to extract them manually.

\begin{figure}\centering
\begin{tabular}{|l||l|l|l|l|l|l|}\hline
         & \multicolumn{2}{c|}{PVS source} & \multicolumn{2}{c|}{PVS $\to$ XML}& \multicolumn{2}{c|}{XML $\to$ \omdoc}\\
         & size/gz  & check time & result size/gz & run time  & result size/gz & run time\\\hline\hline
Prelude  & 189.7/46.6kB & 33s & 23.5/.67MB &  11s  & 83.3/1.6MB & 3m41s \\\hline
NASA Lib & 1.9/.426MB  & 23m25s & 387.2/8.9MB & 3m11s & 2.5/.04GB & 58m56s \\\hline
\end{tabular}
\caption{File Sizes of the PVS Import at Various Stages}\label{fig:sizes}
\end{figure}
%%% Local Variables: 
%%% mode: latex
%%% TeX-master: "paper"
%%% End: 


%  LocalWords:  hline centering lstlisting EquivClos fig:pvsxml fig:pvsscala Expr omdoc
%  LocalWords:  lstinline fig:pvsomdoc pvspi fig:sizes
